%%%%%%%%%%%%%%%%%%%%%%%%%%%%%%%%%%%%%%%%%%%%%%%%%%%%%%%%%%%%%%%%%%%
%%%                                                             %%%
%%%      File: poster1.tex                                      %%%
%%%    Author: Sabine Schulte im Walde                          %%%
%%%   Purpose: Research Seminar                                 %%%
%%%                                                             %%%
%%%%%%%%%%%%%%%%%%%%%%%%%%%%%%%%%%%%%%%%%%%%%%%%%%%%%%%%%%%%%%%%%%%


% a0poster Portrait Poster
% LaTeX Template
% Version 1.0 (22/06/13)
%
% The a0poster class was created by:
% Gerlinde Kettl and Matthias Weiser (tex@kettl.de)
% 
% This template has been downloaded from:
% http://www.LaTeXTemplates.com
%
% License:
% CC BY-NC-SA 3.0 (http://creativecommons.org/licenses/by-nc-sa/3.0/)


%----------------------------------------------------------------------------------------
%	PACKAGES AND OTHER DOCUMENT CONFIGURATIONS
%----------------------------------------------------------------------------------------

\documentclass[a0,portrait]{a0poster}

\usepackage{multicol} % This is so we can have multiple columns of text side-by-side
\columnsep=100pt      % This is the amount of white space between the columns in the poster
\columnseprule=3pt    % This is the thickness of the black line between the columns in the poster

\usepackage[svgnames]{xcolor} % Specify colors by their
                              % 'svgnames'. For a full list of all
                              % colors available see here:
                              % http://www.latextemplates.com/svgnames-colors

\usepackage{times} % Use the times font

\usepackage{graphicx}     % Required for including images
\graphicspath{{figures/}} % Location of the graphics files
\usepackage{booktabs}     % Top and bottom rules for table
\usepackage[font=small,labelfont=bf]{caption}   % Required for specifying captions to tables and figures
\usepackage{amsfonts, amsmath, amsthm, amssymb} % Required for math fonts, symbols and environments
\usepackage{wrapfig}      % Allows wrapping text around tables and figures
%\usepackage{graphicx} % Allows including images
%\usepackage{booktabs} % Allows the use of \toprule, \midrule and \bottomrule in tables
\usepackage{hyperref}
\usepackage{xcolor}
\usepackage{amsmath}


\begin{document}

%----------------------------------------------------------------------------------------
%	POSTER HEADER 
%----------------------------------------------------------------------------------------

% The header is divided into two boxes:
% The first is 75% wide and houses the title, subtitle, names, university/organization and contact information
% The second is 25% wide and houses a logo for your university/organization or a photo of you
% The widths of these boxes can be easily edited to accommodate your content as you see fit

\begin{minipage}{0.75\linewidth}
  \veryHuge \color{NavyBlue} \textbf{Research Seminar at the IMS} \color{Black} % Title
  \vspace{+1cm}\\
  \Huge \textit{Presentation on paper : GQA: A New Dataset for Real-World Visual Reasoning and Compositional Question Answering. Hudson et al. (2019)}\\[2cm] % Subtitle
  \huge \textbf{Shawon Ashraf}\\[0.5cm] % Author(s)
  \huge IMS, University of Stuttgart\\[0.4cm] % University/Organization
  \Large \texttt{shawon.ashraf@ims.uni-stuttgart.de}\\
\end{minipage}
%
\begin{minipage}{0.25\linewidth}
  \includegraphics[width=15cm]{logo_red_text_left}\\
\end{minipage}

\vspace{5cm} % A bit of extra whitespace between the header and the poster content

%----------------------------------------------------------------------------------------

\begin{multicols}{2} % This is how many columns your poster will be broken into. A portrait poster is generally split into 2 columns.


%----------------------------------------------------------------------------------------
%	INTRODUCTION/MOTIVATION
%----------------------------------------------------------------------------------------

\color{SaddleBrown} % SaddleBrown color for the introduction

\section*{Introduction/Motivation}
GQA \cite{hudson2019gqa} is a new dataset for Visual Commonsense Reasoning and Visual Question Answering. It aims to mitigate the issue with previous VQA datasets (such as CLEVR\cite{johnson2017clevr} and VQA\cite{antol2015vqa}) where models can make random guesses about answers and still end up with fairly good accuracy scores. Furthermore, state of the art models can bank on real world priors to exploit the relatively low diverse features in them. GQA uses Scene Graphs \cite{ji2020action} to represent object relationships and  requires the models  to actually understand the visual and linguistic information contained in the data.  
\vspace{+5mm}\\


%----------------------------------------------------------------------------------------
%	OBJECTIVES
%----------------------------------------------------------------------------------------

\color{DarkSlateGray} % DarkSlateGray color for the rest of the content

\vspace{+3cm}
\section*{Main Objectives}

\begin{enumerate}
\item Include inter object relationships in data features so models have to learn to reason about them.
\item Use scene graphs to represent relationships.
\item Evaluate Entailment, Plausability, Validity and Grounding of the predicted answers alongside accuracy.
\item Introduce a baseline and its comparison against human level performance as a reference for future systems.
\end{enumerate}


% --------------------- Scene Graphs
\vspace{+2cm}
\section*{Scene Graphs and Object Relationships}

\begin{center}\vspace{1cm}
  \includegraphics[width=0.4\linewidth]{presentation/gqa_1.png}
  \captionof{figure}{\color{Green} An Example Picture}
\end{center}\vspace{1cm}

\begin{center}\vspace{1cm}
  \includegraphics[width=0.6\linewidth]{presentation/scene_graph.png}
  \captionof{figure}{\color{Green} Scene Graph for the example picture}
\end{center}\vspace{1cm}

\begin{center}\vspace{1cm}
  \includegraphics[width=0.6\linewidth]{poster/figures/example.png}
  \captionof{figure}{\color{Green} Scene Graph}
\end{center}\vspace{1cm}

%----------------------------------------------------------------------------------------
%	MATERIALS AND METHODS
%----------------------------------------------------------------------------------------

%\vspace{+2cm}
\columnbreak
\section*{Methods}

\begin{enumerate}
\item Combining features from the Visual Gnome Dataset with linguistic patterns from VQA 2.0
\item Using semantic features from questions and images to reduce bias.
\item A more structured and normalized orientation for the scene graphs, representing it in terms of object, attributes and relationships.
\item A Questions Engine which combines relationships between objects in images and the linguistic features in descriptions of the images to generate questions.
\item Ut adipiscing accumsan sapien, sit amet pretium.
\item Baseline Models: CNN, LSTM, MAC\cite{hudson2018compositional}, CNN + LSTM
\item Baselines compared against human performance.
\end{enumerate}


%------------------------------------------------

\vspace{+1cm}
\subsection*{Dataset Information}
   \begin{center}\vspace{1cm}
        \begin{tabular}{@{}lll@{}}
            \toprule
            Questions  & Images  & Labels \\ \midrule
            22,669,678 & 113,018 & 1878 \\  \bottomrule
        \end{tabular}
        \caption{Questions, Images and Labels in GQA}
    \end{center}
    

\vspace{+1cm}
\subsection*{Evaluation Settings}
        \begin{itemize}
            \item Test Size: 5000; \textbf{Crowd Sourced}
            \item Train:Validation $\rightarrow$ $88:12$
        \end{itemize}


%----------------------------------------------------------------------------------------
%	RESULTS 
%----------------------------------------------------------------------------------------


%\vspace{+2cm}
\section*{Results}
\begin{center}\vspace{1cm}
  \includegraphics[width=0.8\linewidth]{presentation/gqa_res.png}
  \captionof{figure}{\color{Green} Baseline models compared against human performance}
\end{center}\vspace{1cm}





%


%----------------------------------------------------------------------------------------
%	CONCLUSIONS
%----------------------------------------------------------------------------------------

\color{SaddleBrown} % SaddleBrown color for the conclusions to make them stand out

\vspace{+2cm}
\section*{Conclusions}

\begin{itemize}
\item The GQA dataset provides a new benchmark for testing semantic reasoning in models. 
\item There is a challenge task: \hyperlink{https://cs.stanford.edu/people/dorarad/gqa/challenge.html}{https://cs.stanford.edu/people/dorarad/gqa/challenge.html} on the dataset. The authors provide test sets and a baseline to compare implementations against. 
\end{itemize}

\color{DarkSlateGray} % Set the color back to DarkSlateGray for the rest of the content


%----------------------------------------------------------------------------------------
%	FORTHCOMING RESEARCH
%----------------------------------------------------------------------------------------


%----------------------------------------------------------------------------------------
%	REFERENCES
%----------------------------------------------------------------------------------------

\vspace{+1cm}
\nocite{*} % Print all references regardless of whether they were cited in the poster or not
\bibliographystyle{plain} % Plain referencing style
\bibliography{sample1} % Use the example bibliography file sample.bib


%----------------------------------------------------------------------------------------
%	ACKNOWLEDGEMENTS
%----------------------------------------------------------------------------------------



%----------------------------------------------------------------------------------------

\end{multicols}

\end{document}


%% --- END OF FILE
